\documentclass[../main.tex]{subfiles}
\begin{document}


\section{Overview}

This chapter presents the methodology of the proposed pipeline for automatic personality recognition
from written text using the Big Five (OCEAN) model.

The pipeline is organised as a retrieval-augmented generation (RAG) pipeline that operates
in two separated phases. The first phase, the \emph{index-building phase}, is executed once over the training corpus: each training
essay is processed by a large language model (LLM) to extract a thirty-facet psychological representation of
the writer. This representation is then embedded, fused with an embedding of the raw essay, and stored in a
similarity index. The corresponding thirty numerical facet scores are retained as structured metadata.
The second phase is inference phase executed using test essays. During inference phase, all the test essays are
also profiled using the same process as in the first phase. The index is queried for the most psychologically
similar training essays, the retrieved exemplars are formatted into a prompt,
and an LLM outputs the final \texttt{high}/\texttt{low} prediction for each of the
five traits. The remainder of this chapter describes each stage in detail.

The two-phase architecture is summarised in Figure~\ref{fig:pipeline}.

\begin{figure}[htbp]
  \centering
  \includegraphics[width=0.95\textwidth]{Figure/pipeline.png}
  \caption{Overview of the two-phase RAG pipeline for automatic personality recognition.
           The index-building phase (top) processes training essays once; the inference
           phase (bottom) queries the index and generates trait predictions.}
  \label{fig:pipeline}
\end{figure}


\section{The Thirty-Facet Psychological Profiler}
\label{sec:profiler}

\subsection{Motivation}
\label{subsec:profiler-motivation}
A straightforward approach for text-based personality recognition is to embed 
raw essays using a general-purpose encoder and retrieve similar samples for a target 
text based on vector distance. However, this approach has two notable limitations.

First, general-purpose sentence encoders place two essays close together when they 
share vocabulary and topic. This is problematic because personality detection relies 
less on surface-level content and more on psychological signals such as emotional 
expression, cognitive style, and references to social relationships. In other words, 
two people may write about the same topic in very different ways depending on their 
personality, yet a general-purpose encoder would treat their essays as similar. 
This gap between topical similarity and psychological similarity can be thought of as 
a surface confound — where the retrieval mechanism returns essays that read alike 
rather than essays written by psychologically similar authors. Consequently, the 
examples provided to the large language model reinforce this confound rather than 
counteracting it.

Second, even when relevant essays are retrieved, providing the model with full-length 
essays as in-context examples may introduce noise and dilute the personality-relevant 
signals within a long body of text, making it harder for the model to focus on the 
features that matter.

\subsection{The NEO-PI-R Facet Model}
\label{subsec:neo-pi-r}
The schema used by the profiler is the facet structure of the Revised NEO
Personality Inventory (NEO-PI-R) \cite{costa2008revised}. The
NEO-PI-R decomposes each of the five broad OCEAN traits --- Openness to
Experience, Conscientiousness, Extraversion, Agreeableness and Neuroticism ---
into six narrower \emph{facets}, yielding a total of thirty facets. For
example, Extraversion is decomposed into \emph{warmth}, \emph{gregariousness},
\emph{assertiveness}, \emph{activity}, \emph{excitement-seeking} and
\emph{positive emotion}. Facets are valuable for two reasons. They are
psychologically interpretable: a reader can examine a facet profile and
understand it. And they are more discriminative than the broad trait alone: two
writers may both be moderately extraverted, yet one is gregarious-but-passive
and the other assertive-but-solitary, a distinction that the six Extraversion
facets capture but a single Extraversion score discards.
The complete catalogue of thirty facets used by the system is listed in
Table~\ref{tab:facets}.

\begin{longtable}{p{1.2cm}p{2.8cm}p{3.6cm}p{5.4cm}}
  \caption{The thirty NEO-PI-R facets used by the psychological profiler.
  Each facet receives one independent signal
  in $\{\textit{high}, \textit{mod}, \textit{low}, \textit{none},
  \textit{n/e}\}$.}
  \label{tab:facets}\\
  \toprule
  \textbf{Code} & \textbf{Facet} & \textbf{Parent trait} & \textbf{Gloss} \\
  \midrule
  \endfirsthead
  \multicolumn{4}{c}{\textit{Table~\ref{tab:facets} continued from previous page}}\\
  \toprule
  \textbf{Code} & \textbf{Facet} & \textbf{Parent trait} & \textbf{Gloss} \\
  \midrule
  \endhead
  \midrule \multicolumn{4}{r}{\textit{continued on next page}}\\
  \endfoot
  \bottomrule
  \endlastfoot
  N1 & anxiety           & Neuroticism       & free-floating worry, fear, tension \\
  N2 & hostility         & Neuroticism       & anger, irritation, resentment toward others \\
  N3 & depression        & Neuroticism       & sadness, hopelessness, low mood \\
  N4 & self-conscious    & Neuroticism       & shame, embarrassment, social discomfort \\
  N5 & impulsiveness     & Neuroticism       & difficulty resisting urges or delaying \\
  N6 & vulnerability     & Neuroticism       & feeling overwhelmed, helpless under stress \\
  E1 & warmth            & Extraversion      & affectionate, friendly engagement with others \\
  E2 & gregariousness    & Extraversion      & preference for company, group settings \\
  E3 & assertiveness     & Extraversion      & leadership, dominance, decisive voice \\
  E4 & activity          & Extraversion      & energetic pace, busy schedule, motion \\
  E5 & excitement-seek   & Extraversion      & thrill-seeking, novelty appetite \\
  E6 & positive emotion  & Extraversion      & joy, enthusiasm, optimism \\
  O1 & fantasy           & Openness          & vivid imagination, daydreaming \\
  O2 & aesthetics        & Openness          & appreciation of art, beauty, nature \\
  O3 & feelings          & Openness          & emotional depth, introspective awareness \\
  O4 & actions/variety   & Openness          & openness to new experiences, change \\
  O5 & ideas             & Openness          & intellectual curiosity, abstract thought \\
  O6 & values            & Openness          & willingness to re-examine values, norms \\
  A1 & trust             & Agreeableness     & belief in others' good intentions \\
  A2 & straightforward   & Agreeableness     & candor, frankness, low manipulation \\
  A3 & altruism          & Agreeableness     & active concern for others' welfare \\
  A4 & compliance        & Agreeableness     & deference, low antagonism in conflict \\
  A5 & modesty           & Agreeableness     & humility, low self-promotion \\
  A6 & tender-minded     & Agreeableness     & sympathy, empathy with others' pain \\
  C1 & competence        & Conscientiousness & sense of capability, preparedness \\
  C2 & order             & Conscientiousness & tidiness, organization, structure \\
  C3 & dutifulness       & Conscientiousness & adherence to obligations, ethics \\
  C4 & achievement       & Conscientiousness & striving for goals, ambition \\
  C5 & self-discipline   & Conscientiousness & follow-through, resisting distraction \\
  C6 & deliberation      & Conscientiousness & careful planning, thinking before acting \\
\end{longtable}

\subsection{The Controlled Signal Vocabulary}
\label{subsec:signal-vocab}

For every one of the thirty facets the profiler must emit exactly one
\emph{signal} drawn from a small controlled vocabulary. In the code this
vocabulary is the constant \texttt{SIGNAL\_VOCAB} and contains five values:
\texttt{high}, \texttt{mod}, \texttt{low}, \texttt{none} and \texttt{n/e}. The
first three are the obvious ordinal levels --- clear evidence in the high
direction, mixed or weak evidence, and clear evidence in the low direction
respectively. The remaining two values, \texttt{none} and \texttt{n/e}, deserve
careful explanation because the distinction between them is a deliberate design
decision.

The value \texttt{none} means that the facet \emph{is} observable in this genre
of writing but the particular writer expresses it neutrally or not at all. The
value \texttt{n/e}, read ``not evidenced'', means something fundamentally
different: the facet is not realistically observable in solo, introspective
journal writing at all. Without an explicit
``not-evidenced'' value, a profiler forced to choose among
\texttt{high}/\texttt{mod}/\texttt{low} would be pushed to fabricate evidence
on exactly those facets. Making \texttt{n/e} a first-class member of the
vocabulary is intended to reduce that pressure. As
Section~\ref{subsec:facet-encoding} shows, \texttt{none} and \texttt{n/e} are
deliberately collapsed to the same numerical value of zero, to avoid treating
\emph{absence of evidence} as \emph{evidence of a low score}.

\subsection{The Linguistic Fingerprint}
\label{subsec:linguistic}

In addition to the thirty facet lines, the profiler emits a short
\emph{linguistic fingerprint}: seven fixed keys --- \emph{pronouns},
\emph{emotion}, \emph{tense}, \emph{cognitive}, \emph{social}, \emph{register}
and \emph{length} --- each with a one-line free-text value. These keys follow
the strongest meta-analytic predictors of personality in the LIWC and
open-vocabulary literature~\cite{Pennebaker2007TheDA}.
They capture genre-level and stylistic information that the facet schema
deliberately suppresses, and they are carried into the embedded profile text
where they contribute to retrieval similarity.

% =============================================================================
\section{Profile Embedding and Trait-Aligned Retrieval}
\label{sec:retrieval}
% =============================================================================

\subsection{Profile Embedding}
\label{subsec:profile-embedding}

Once the profiler produces a structured profile, that profile is serialised to
text and passed through a sentence encoder to produce the retrieval vector. By
embedding the \emph{profile} rather than the raw essay, the retrieval space is
defined over the psychological facet vocabulary rather than surface lexis.

\subsection{Facet Vector Encoding}
\label{subsec:facet-encoding}

In parallel with the profile embedding, each profile is encoded as a point
in $\mathbb{R}^{30}$, one coordinate per facet. The categorical signal is
mapped to an ordinal value:
%
\begin{equation}
  \label{eq:signal-map}
  \texttt{high}\mapsto +1.0,\quad
  \texttt{mod}\mapsto +0.5,\quad
  \texttt{none}\mapsto 0.0,\quad
  \texttt{n/e}\mapsto 0.0,\quad
  \texttt{low}\mapsto -1.0.
\end{equation}
%
The scale is signed and symmetric; \texttt{mod} contributes only half the
weight of a confident signal; \texttt{none} and \texttt{n/e} both map to
zero so that absence of evidence neither attracts nor repels a candidate.
The resulting raw vector is L2-normalised:
%
\begin{equation}
  \label{eq:facet-vector}
  \boldsymbol{\phi}(p) \;=\; \frac{\mathbf{u}(p)}{\lVert \mathbf{u}(p) \rVert_2}.
\end{equation}

\subsection{Trait-Masked Facet Similarity}
\label{subsec:trait-slice}

When predicting a specific trait $\tau$, only the six facets of that trait are
relevant; the remaining twenty-four introduce cross-trait noise. Let
$S_\tau \subset \{1,\dots,30\}$ denote the six column indices belonging to
$\tau$. The trait-masked similarity between query $q$ and candidate $i$ is
%
\begin{equation}
  \label{eq:facet-cosine}
  \operatorname{sim}^{\mathrm{facet}}_{\tau}(q, i)
  \;=\;
  \cos\!\big(
    \boldsymbol{\phi}_q[S_\tau],\;
    \boldsymbol{\phi}_i[S_\tau]
  \big),
\end{equation}
%
where both sub-vectors are re-normalised after masking to remove any magnitude
bias from the discarded facets.

\subsection{Hybrid Retrieval Score}
\label{subsec:hybrid-score}

The retriever combines the dense profile-embedding similarity with the
trait-masked facet similarity into a single score. For query $q$, candidate
$d$, and trait $\tau$:
%
\begin{equation}
  \label{eq:hybrid}
  \operatorname{score}_\tau(q, d)
  \;=\;
  \beta_\tau \cdot \operatorname{sim}^{\mathrm{dense}}(q, d)
  \;+\;
  \gamma_\tau \cdot \operatorname{sim}^{\mathrm{facet}}_{\tau}(q, d).
\end{equation}
%
The weights $\beta_\tau$ and $\gamma_\tau$ are per-trait, with $\beta_\tau = 1 - \gamma_\tau$,
and are selected independently for each trait by maximising mean-match-rate on the
validation set (see Chapter~\ref{chap:experiments}). Traits highly susceptible to
genre confound --- Conscientiousness and Neuroticism --- receive high $\gamma_\tau$,
placing more weight on the structured facet channel; traits where surface text
is more diagnostic may receive lower $\gamma_\tau$.

Retrieval proceeds in two steps: an approximate nearest-neighbour search over
the profile-embedding index returns a candidate pool, which is then re-ranked
by the hybrid score to yield the final top-$k$ exemplars.

\subsection{Full-Profile Exemplar Rendering}
\label{subsec:exemplar-rendering}

Before being inserted into the prompt, each retrieved exemplar is rendered as
a \emph{full-profile}: all thirty facets are included, grouped by trait
(N\,/\,E\,/\,O\,/\,A\,/\,C), together with their evidence strings and a
linguistic summary block. The classifier therefore receives the complete
psychological portrait of each exemplar --- not raw essay text ---
alongside its known label.


% =============================================================================
\section{The Reasoned Inference Stage}
\label{sec:prompt}
% =============================================================================

\subsection{Structured Reasoning Format}

The language model is required to emit its output as six XML blocks in a fixed
order: \texttt{<evidence>}, \texttt{<facet\_check>}, \texttt{<cue\_tally>},
\texttt{<ex\-am\-ple\_a\-lign\-ment>}, \texttt{<verdict>}, and \texttt{<label>}.
No text is permitted outside these tags.

The \texttt{<verdict>} must be consistent with the \texttt{<cue\_tally>},
coupling the final answer to the audited evidence count rather than to the
tone of retrieved exemplars.

\subsection{Prompt Structure}

The user template provides the trait name, the HIGH and LOW definitions of the
trait, the retrieved trait-sliced exemplars, and the test essay. The template
is shown in Listing~\ref{lst:reasoned}.

\begin{lstlisting}[style=promptbox,
  caption={Reasoned RAG user template.},
  label={lst:reasoned}]
Trait: {trait_name}

HIGH {trait_name}: {definition_high}
LOW  {trait_name}: {definition_low}
---

The following are the {top_k} most similar texts from the training set,
with their known labels and extracted psychological evidence. Use them
as calibration anchors - note how each one's cues map to its label.

{similar_context}

---

Text to classify:
<text>

Reason step by step and emit your output in the EXACT XML structure
specified in the system message.
\end{lstlisting}

Retrieval is performed per trait: a separate hybrid query is issued for each
of the five traits so that the masked facet score of
Equation~\eqref{eq:facet-cosine} is specialised to the trait being predicted.
The default exemplar count is $k = 5$ per trait; the effect of $k$ is examined
in Chapter~\ref{chap:experiments}.

% =============================================================================
\section{Chapter Summary}
\label{sec:methodology-summary}
% =============================================================================

This chapter has described the three components of the proposed pipeline.
The thirty-facet psychological profiler maps raw essay text to a structured
NEO-PI-R facet profile using a controlled signal vocabulary, replacing
surface-style representations with a psychologically grounded intermediate
form. The hybrid retrieval mechanism encodes each profile as both a dense
embedding and a 30-dimensional ordinal facet vector, and retrieves the
$k=5$ most psychologically similar training exemplars via a per-trait
weighted combination of dense and structured similarity. The structured
reasoning stage presents those exemplars as calibration anchors and
enforces a six-block XML output format that separates stable trait cues
from transient state cues before issuing the final verdict. The
experimental evaluation of this pipeline is presented in
Chapter~\ref{chap:experiments}.


\end{document}
